\documentclass[11pt]{article}
\usepackage[a4paper, top=18mm, bottom=18mm, left=20mm, right=20mm]{geometry}
\usepackage[T2A]{fontenc}
\usepackage[utf8]{inputenc}
\usepackage[ukrainian]{babel}
\usepackage{graphicx}
\setlength{\parindent}{0pt}
\setlength{\parskip}{4pt}
\pagestyle{empty}
\begin{document}
%begin
{\large\textbf{final-open-40}}

\textbf{Варіант \No\,1}\hfill \textit{ПІБ:}\,\underline{\hspace{60mm}}
\vspace{4mm}

%beginitem
1. Що буде виведено за виконання фрагменту коду?
Записати ВИРАЗ, значенням якого є словник, ключами якого є цілі числа від 13 до 2003 (включно), що не діляться на 5, а ключам відповідають їх квадрати 

%enditem
%beginitem
2. Що буде виведено за виконання фрагменту коду?

%code
def f(a,b):
    c=72
    if a>-2:
        return 2
    if b>=4:
         c=9
    else: 
        c=0
    return c

print(f(-6,-6))
%code

%enditem
%beginitem
3. Що буде виведено за виконання фрагменту коду?

%code
#include <iostream>
using namespace std;

bool f(int n){
    cout<<"f";
    return n;
}

int main(){
    cout<<(f(-6) && f(0));
    return 0;
}
%code

%enditem
%beginitem
4. Що буде виведено за виконання фрагменту коду?
Рядки журналу через = містять ім'я користувача (ідентифікатор), час події,тип події, код події, наприклад
\begin{verbatim}
s="username = 2022-05-05 13:26 = ERROR = 404"
\end{verbatim}
у коді 
\begin{verbatim}
res = re.fullmatch('???', s) 
s = ''
code_str = ??? 
\end{verbatim}
чим треба замінити кожен з ???, щоб значенням code\_str був код події? 



%enditem
%beginitem
5. Що буде виведено за виконання фрагменту коду?

%code
a,b,c=2,9,0
def f(a):
    a=4
    b=3
    c=5
    return a+b+c

a,b,c=5,4,6
print(f(a),a,b,c)
%code

%enditem
%beginitem
6. Що буде виведено за виконання фрагменту коду?

%code
#include <iostream>
using namespace std;

int f(int a, int b){
    int c = 72;
    if (a > -2)
        return 2;
    if (b >= 4)
         c = 9;
    else 
        c = 0;
    return c;
}

int main(){
    cout << f(-6, -6);
    return 0;
}
%code

%enditem
%beginitem
7. Що буде виведено за виконання фрагменту коду?
%code

try:
    print(2,end="")
    print(int("9"),end="")
    print(0,end="")
except ValueError: print(5,end="")
except TypeError: print(4,end="")
else: print(6,end="")
finally: print(7,end="")

%code

%enditem
%beginitem
8. Що буде виведено за виконання фрагменту коду?
True>"4"
%enditem
%beginitem
9. Що буде виведено за виконання фрагменту коду?

print('R', end=' ')
for c in range(-12,-12,2):
    if c>5:
        continue
        print(c, end=' ')
    if c>5:
        break
    else:
        print(c, end=' ')
    print(c)

%enditem
%beginitem
10. Що буде виведено за виконання фрагменту коду?

%code
print('R', end=' ')
c=[10,11,12,13,14]
c[-7:-7]='11'
print(c)
%code


%enditem
%beginitem
11. Що буде виведено за виконання фрагменту коду?
a,b,c=2,9,0
   def f(a):
     global c     a=4
     b+=3
     c=5
     return a+b+c
   a,b,c=5,4,6
   print(f(a),a,b,c)
%enditem
%beginitem
12. Що буде виведено за виконання фрагменту коду?

%code
print(True>"4")
%code

%enditem
%beginitem
13. Що буде виведено за виконання фрагменту коду?
Змінні $next$ та $prev$ вузла списку позначають наступний та попередній за ним вузол відповідно. Починаючи з голови, у список записано послідовні цілі числа від $-68$ до $17$. Нехай змінна $p$ позначає вузол, що зберігає значення $-6$.
Яке значення зберігається у вузлі $p.prev.prev.prev.prev.prev$ ?

%enditem
%beginitem
14. Що буде виведено за виконання фрагменту коду?
k=8
   while k>2: k+=-3
   print(k)
%enditem
%beginitem
15. Що буде виведено за виконання фрагменту коду?
c=[10,11,12,13,14]; c[-7:-7]='11'; print(c)
%enditem
%beginitem
16. Що буде виведено за виконання фрагменту коду?
try:
    k = 8
    while k>2:
        k += -3
    print(k, end=';')
except:
    print('ERR')

%enditem
%beginitem
17. Що буде виведено за виконання фрагменту коду?
for c in range(-12,-12,2):
     if c>5: continue
     print(c,end=' ')
     if c>5: break
   else: print(c,end=' ')
   print(c)
%enditem
%beginitem
18. Що буде виведено за виконання фрагменту коду?
Записати ВИРАЗ, значенням якого є список, що складається з квадратівцілих чисел від 13 до 2003 (включно), що не діляться на 5, записаних в обернено\-му порядку.

%enditem
%beginitem
19. Що буде виведено за виконання фрагменту коду?
Записати ВИРАЗ, значенням якого є список, що складається з цілих чисел від 13 до 2003 (включно), причому спочатку за зростанням записано всі, що не діляться на 5, а потім решту за зростанням.

%enditem
%beginitem
20. Що буде виведено за виконання фрагменту коду?

%code
a,b,c=2,9,0
def f(a):
global c    a=4
    b=3
    c=5
    return a+b+c

a,b,c=5,4,6
print(f(a),a,b,c)
%code

%enditem
%beginitem
21. Що буде виведено за виконання фрагменту коду?
%code
for c in range(-12,-12,2):
    if c>5:
        continue
        print(c,end=' ')
    if c>5:
        break
else:
    print(c, end=' ')
print(c, end=' ')
%code

%enditem
%beginitem
22. Що буде виведено за виконання фрагменту коду?

%code
print('R', end=' ')
k=2
while k<8:
    k+=1
print(k)
%code


%enditem
%beginitem
23. Що буде виведено за виконання фрагменту коду?
a = 2
b = 9
c = 0

def f(a):
global c    a = 4
    b += 3
    c = 5
    return a + b + c

a = 5
b = 4
c = 6

try:
    print(f(a), a, b, c, sep=';')
except:
    print('ERR', a, b, c, sep=';')

%enditem
%beginitem
24. Що буде виведено за виконання фрагменту коду?
try:
    for c in range(-12, -12, 2):
        if c>5:
            continue
        print(c, end=';')
        if c>5:
            break
    else:
        print(c,end=';')
    print(c)
except:
    print('ERR')

%enditem
%beginitem
25. Що буде виведено за виконання фрагменту коду?
%code

x,y,z=2,3,0
z,x,x,z=x,8,z,y
print(x,y,z)
%code

%enditem
%beginitem
26. Що буде виведено за виконання фрагменту коду?

%code
True>"4"
%code

%enditem
%beginitem
27. Що буде виведено за виконання фрагменту коду?
-2**-2+1
%enditem
%beginitem
28. Що буде виведено за виконання фрагменту коду?

%code
print('R', end=' ')
c = ('0','1','2','3','4','5','6','7','8','9')
print(c[-12:-12:2])
%code


%enditem
%beginitem
29. Що буде виведено за виконання фрагменту коду?
c=[10,11,12,13,14]; c[-7:-7]='11'; print(c)
%enditem
%beginitem
30. Що буде виведено за виконання фрагменту коду?
Записати ВИРАЗ, значенням якого є список, що складається з цілих чисел від 13 до 2003 (включно), причому спочатку за зростанням записано всі, що не діляться на 5, а потім решту за зростанням.

%enditem
%beginitem
31. Що буде виведено за виконання фрагменту коду?
Записати ВИРАЗ, значенням якого є список, що складається з квадратівцілих чисел від 13 до 2003 (включно), що не діляться на 5, записаних в оберненому порядку.
%enditem
%beginitem
32. Що буде виведено за виконання фрагменту коду?

%code
True>"4"
%code

%enditem
%beginitem
33. Що буде виведено за виконання фрагменту коду?
k=2
   while k<8: k+=1
   print(k)
%enditem
%beginitem
34. Що буде виведено за виконання фрагменту коду?

%code
6/6//6*6
%code

%enditem
%beginitem
35. Що буде виведено за виконання фрагменту коду?

%code
a,b,c=2,9,0
def f(a):
global c    a=4
    b+=3
    c=5
    return a+b+c

a,b,c=5,4,6
print(f(a),a,b,c)
%code

%enditem
%beginitem
36. Що буде виведено за виконання фрагменту коду?
try:
    for c in range(-12, -12, 2):
        if c>5:
            continue
        print(c, end=';')
        if c>5:
            break
    else:
        print(c,end=';')
    print(c)
except:
    print('Z')

%enditem
%beginitem
37. Що буде виведено за виконання фрагменту коду?
%code
for c in range(-12,-12,2):
    if c>5:
        continue
        print(c,end=' ')
    if c>5:
        break
else:
    print(c, end=' ')
print(c, end=' ')
%code

%enditem
%beginitem
38. Що буде виведено за виконання фрагменту коду?

%code
print(4 >= 4.0 > 5 == True)
%code

%enditem
%beginitem
39. Що буде виведено за виконання фрагменту коду?
%code

a,b,c=8,9,6
def f(a,b=7,c):
    print(a,b,c,end="")

f(2,a=3)
print(a,b,c)

%code

%enditem
%beginitem
40. Що буде виведено за виконання фрагменту коду?
try:
    k = 2
    while k<8:
        k += 1
    print(k, end=';')
except:
    print('ERR')

%enditem








































%end
\newpage
\end{document}

